\documentclass[11pt]{article}

\usepackage[margin=1in]{geometry}
\usepackage{amsmath, amssymb, amsthm}
\usepackage{mathtools}
\usepackage{booktabs}
\usepackage{tabularx}
\usepackage{graphicx}
\graphicspath{{figures/}{../figures/}}
\usepackage[hidelinks]{hyperref}
\usepackage[authoryear,round]{natbib}
\bibpunct{(}{)}{,}{a}{}{,}

\newcommand{\R}{\mathbb{R}}
\newcommand{\E}{\mathbb{E}}
\newcommand{\F}{\mathcal{F}}
\newcommand{\pluspart}[1]{\left[#1\right]_+}

\theoremstyle{plain}
\newtheorem{theorem}{Theorem}
\newtheorem{lemma}{Lemma}
\newtheorem{proposition}{Proposition}
\theoremstyle{definition}
\newtheorem{definition}{Definition}
\newtheorem{assumption}{Assumption}
\theoremstyle{remark}
\newtheorem{remark}{Remark}

\title{Optimal Trading Filters:\\ a Wiener--Hopf Approach}
\author{Matteo Smerlak\thanks{Capital Fund Management, 23 rue de l'Universit\'e, 75007 Paris, France.\ \texttt{matteo.smerlak@cfm.com}}}
\date{}

\begin{document}
\maketitle


\begin{abstract}
A trader holding a predictive signal must weigh the expected gain against the cost of acquiring the position and the risk of holding it. Under the propagator model of market impact, trading-induced price changes involve the entire history of past trades. Existing treatments of the resulting path-wise quadratic program characterize the optimum implicitly, as the solution of an integral equation or an equivalent forward--backward system. In this paper I show how stationary solutions can be computed analytically from the signal and its forecasts using a Wiener-Hopf factorization of the friction operator. For rational frictions --- inventory risk, temporary cost, and exponential resilience --- the trading filter reduces to finitely many exponential moving averages, and we recover the Neuman--Vo\ss{} solution away from horizon boundaries. In the (empirically more realistic) case of scale-free impact kernels, the optimal trading filter decays algebraically rather than exponentially at large frequencies; in the limit of small risk aversion, the target position is a fractional integral of the signal. Unlike its exponential (G\^arleanu--Pedersen) cousin, this solution never surfs its own impact or generates block trades.
\end{abstract}

% =====================================================================
\section{Introduction}\label{sec:intro}

\subsection{The gain--risk--cost problem}\label{sec:problem}

Optimal execution studies how to build or unwind a position when trading moves the price against the trader \citep{AlmgrenChriss2001,GatheralSchiedSlynko2012}. A return-predicting signal sharpens the problem: a forecast invites a position, but the trade that captures it pushes the price away, commits capital to risk, and takes time the forecast may not survive. The trader weighs the gain a signal promises against the cost of realizing it and the risk of holding it, acting on the information in hand and never on the signal's future.

Concretely, a trader observes a return signal $\mu_t$, the conditional expectation at time $t$ of the asset's instantaneous price drift, and holds a position $x_t$, traded at the rate $u_t=\dot x_t$. Trading incurs friction in three ways. An inventory cost penalizes holding large positions. A temporary impact cost, quadratic in the rate, prices the demand for immediate liquidity. A transient impact cost carries the memory that each trade moves the price against later trades, decaying through a propagator kernel $g$ \citep{BouchaudGefenPottersWyart2004,Gatheral2010}.

\citet{AlmgrenChriss2001} split the cost of trading in two: a temporary part, felt only by the trade that causes it, and a permanent part, linear in cumulative volume and left behind for every later trade. Minimizing the resulting mean--variance functional produced closed-form execution schedules. Empirical work then dissolved the sharp split: \citet{LilloFarmerMantegna2003} found impact to be a concave function of size, and \citet{BouchaudGefenPottersWyart2004} found it \emph{transient}, a trade displacing the price and the displacement relaxing, so that permanent and temporary costs are the slow and fast limits of one decaying response. This is the propagator picture, in which the price carries a weighted memory of past trades through a kernel $g$, and the shape of $g$ is the modeling choice.

The tractable shape is \emph{exponential}: \citet{ObizhaevaWang2013} derived $g(t)=e^{-\kappa t}$ from a limit order book of finite depth refilling at a fixed rate, a single resilience time $1/\kappa$ whose fast and slow limits return the temporary and permanent costs of \citet{AlmgrenChriss2001}, and it underlies much of the closed-form execution literature. Transaction data across markets then made clear that the decay is instead a \emph{power law}, $g(t)=t^{-\beta}$ with $\beta\in(0.2,0.6)$ \citep{BouchaudGefenPottersWyart2004,JusselinRosenbaum2020}, scale-free and long-memory; \citet{Gatheral2010} showed a decay of about this form to be what no-dynamic-arbitrage allows, and \citet{JusselinRosenbaum2020} tied the exponent to rough volatility. The exponential kernel offers finite memory and closed forms; the power law has long memory and the empirical evidence.

The trader maximizes the expected gain of the position against the return, net of these costs,
\begin{equation}\label{eq:objective}
\max_{x\ \text{adapted}}\ \E\!\int x_t\mu_t\,dt
\;-\; \frac{\eta}{2}\,\E\!\int \dot x_t^2\,dt
\;-\; \frac{\gamma}{2}\,\E\!\iint g(|t-s|)\,\dot x_t \dot x_s\,dt\,ds
\;-\; \frac{\lambda}{2}\,\E\!\int x_t^2\,dt ,
\end{equation}
over positions adapted to the flow of information: $x_t$ may depend on the signal observed up to time $t$, not on its future. The gain pairs the position against the expected return, the mean--variance objective of portfolio choice; the three costs are the frictions of execution. The integrals run over the whole line; Section~\ref{sec:boundary} poses the same problem on a finite window $[0,T]$ with a terminal inventory constraint. This paper gives exact solutions to \eqref{eq:objective} under both exponential and power-law impact kernels.

\subsection{The adaptedness constraint}\label{sec:adaptedness}

The objective \eqref{eq:objective} pairs a linear gain $\E\int x_t\mu_t\,dt$ against a positive quadratic cost $\tfrac12\E\langle x,Nx\rangle$ and thus evaluates the Legendre transform of the cost at the signal, under the adaptedness constraint. Referred to the position, the three costs define a single \emph{friction operator} $N=-\eta\,\partial_t^2-\gamma\,\partial_t^2(g\ast)+\lambda I$, with Fourier symbol
\begin{equation}\label{eq:symbol}
\hat n(\omega) \;=\; \eta\,\omega^2 \;+\; \gamma\,\hat g(\omega)\,\omega^2 \;+\; \lambda,
\qquad \hat g(\omega) = \int e^{i\omega t}g(|t|)\,dt \ge 0.
\end{equation}
Referred to the trading rate the friction is $Q=N\,(-\partial_t^2)^{-1}$, with symbol $\hat q=\hat n/\omega^2=\eta+\gamma\hat g+\lambda/\omega^2$, the form used for the scale-free kernel of Section~\ref{sec:fractional}. Could the trader see the whole signal path, the first-order condition of \eqref{eq:objective} would be a single operator inversion, the position
\begin{equation}\label{eq:markowitz-anticip}
x \;=\; N^{-1}\mu,
\end{equation}
the continuous-time, memory-laden analogue of the Markowitz rule $w=\lambda^{-1}\Sigma^{-1}\mu$ (weights $w$, covariance $\Sigma$, expected returns $\mu$), with the friction operator in place of the covariance. But the signal path is revealed only as time passes, and \eqref{eq:markowitz-anticip} is acausal.

With inventory cost alone ($\eta=\gamma=0$) the quadratic form is diagonal in time: the first-order condition at time $t$ involves $x_t$ and $\mu_t$ only, and the adapted optimum is the Markowitz rule $x_t=\mu_t/\lambda$ (Section~\ref{sec:recovery}). A temporary cost adds the local operator $-\eta\,\partial_t^2$, still an ordinary differential equation. Transient impact removes this locality: its symbol $\gamma\hat g(\omega)\,\omega^2$ is not a polynomial, so $Nx$ at time $t$ depends on the position over the past \emph{and the future}, and adaptedness --- the requirement that $x_t$ be measurable with respect to $\F_t$ --- becomes binding. Varying \eqref{eq:objective} over adapted perturbations gives the \emph{projected} first-order condition,
\begin{equation}\label{eq:foc}
\E_t\bigl[(Nx^\star)(t)\bigr] \;=\; \mu_t
\qquad\Longleftrightarrow\qquad
P_+\,N\,P_+\,x^\star = \mu,
\qquad (P_+X)_s = \E_s[X_s],
\end{equation}
with $P_+$ the optional projection onto adapted processes. Solving \eqref{eq:foc} requires inverting $P_+NP_+$, and the inverse of a projected operator is not the projection of the inverse, precisely because $N$ is non-local.

\subsection{The Wiener--Hopf method}\label{sec:thesis}

\citet{WienerHopf1931} introduced the technique that resolves this class of problems, originally for convolution equations on a half-line, $\int_0^\infty k(t-s)f(s)\,ds=g(t)$ for $t\ge0$ \citep{Krein1962,Noble1958}. Fourier analysis diagonalizes convolutions on the whole line; a half-line support constraint couples the frequencies, in the same way the filtration constraint couples them in \eqref{eq:foc}. The method restores solvability by factoring the symbol,
\[
\hat k(\omega) \;=\; \hat k_-(\omega)\,\hat k_+(\omega),
\]
with $\hat k_+$ analytic and zero-free in the upper half-plane and $\hat k_-$ in the lower. Analyticity in a half-plane encodes causality: by Paley--Wiener, $\hat k_+$ is the transform of a kernel supported on $t\ge0$ and $\hat k_-$ of one supported on $t\le0$. Operators with causal kernels are triangular with respect to the time order, and triangular operators invert forward in time, staying triangular. The factorization thus splits a two-sided convolution into two one-sided operators, each compatible with the boundary, between which the half-line projection can be threaded. The finite-dimensional analogue is the Cholesky factorization of a positive matrix into triangular factors, computed by Gram--Schmidt along the coordinate order.

The method's classical uses are prediction and control. For linear prediction of a stationary time series the factored symbol is the spectral density, $S=|\hat\psi|^2$, and the causal factor builds the innovations representation \citep{Wiener1949,Whittle1963}. In the trading problem the factored symbol is the \emph{friction} $N$, and the causal factor builds the optimal trade: prediction whitens by the signal's own factor to forecast it; trading whitens by the friction's factor to monetize it. Existence of such a factorization for a general positive operator was established later --- \citet{GohbergKrein1970} for Volterra operators on an interval, \citet{Arveson1975} for positive operators on a nest of subspaces --- and these existence theorems do not by themselves produce the factor. It is explicit in the stationary case, where Fourier diagonalizes $N$ and the causal factor is the scalar Szeg\H{o} outer function; frequency-domain factorization of this kind also solves linear--quadratic dynamic optimization in economics \citep{Whiteman1985}. On a finite horizon translation invariance is lost and the factor is explicit only for the special kernels of Section~\ref{sec:boundary}.

Section~\ref{sec:interior} solves the stationary problem on the whole line in three steps --- factor the friction, predict the signal from its past, and combine the two --- through a projected-inverse identity (Lemma~\ref{lem:pi}) that assembles the adapted optimum for a general signal (Theorem~\ref{thm:general}) and, for stationary Gaussian signals, an explicit trading filter (Theorem~\ref{thm:filter}). Section~\ref{sec:boundary} carries the factorization to a finite horizon, where the Szeg\H{o} factor becomes a terminal-anchored Gohberg--Krein factor and the interior relaxes to the stationary filter. The two remaining sections work out the kernel families: Section~\ref{sec:recovery} the exponential and rational frictions, whose rational symbol gives filters with finitely many moving averages and recovers the classical portfolio and execution rules, and Section~\ref{sec:fractional} the scale-free power law, whose fractional policy, cost of causality $\sin(\pi\beta/2)$, and long-memory liquidation profiles have no rational analogue.

\subsection{Relation to earlier work}\label{sec:relatedwork}

The execution literature has treated pieces of \eqref{eq:objective} by several routes: deterministic liquidation under general propagators \citep{GatheralSchiedSlynko2012}; signal-adaptive trading, introduced for Almgren--Chriss costs by \citet{CarteaJaimungal2013} and developed under temporary and exponential-resilience costs \citep{LehalleNeuman2019,NeumanVoss2022,BankSonerVoss2017}; general propagators through infinite-dimensional stochastic control \citep{AbiJaberNeuman2025} and operator-resolvent methods \citep{AbiJaberNeumanTuschmann2024}; and power-law resilience with Gaussian signals on a finite horizon \citep{FordeSanchezSmith2022}. Each computes an inverse of the same friction operator. The finite-horizon methods produce that inverse as a resolvent series, a Wiener-chaos expansion, or the solution of a Riccati system; the factorization sums those series in closed form when translation invariance is available. Path-dependent and data-driven strategies --- signature methods \citep{KalsiLyonsPerezArribas2020} and reinforcement learning --- approximate the problem and lie outside the closed-form scope here.

The two representations differ in what they expose. Dynamic programming and the backward-SDE methods return a \emph{feedback law}: the trade is a function of the current state --- the position, the outstanding transient-impact profile, and the signal --- with gains fixed by a Riccati or backward SDE, and its dependence on the signal is implicit, carried by the state dynamics one propagates forward. The factorization returns a \emph{filter}: the trade is an explicit linear operator on the signal, a transfer function mapping the observed signal to the trade in closed form --- a gain and phase at each frequency, a fractional derivative for power-law impact, a few moving averages for a rational friction.

\citet{AbiJaberNeuman2025} carry the feedback route to general propagators, characterizing the optimum through a free-boundary backward SDE and an operator Riccati equation, with existence and uniqueness for any Volterra propagator and any adapted signal and closed forms recovered by kernel-specific reductions; the variational form of the same problem is a stochastic Fredholm equation, extended to linear constraints by Lagrange multipliers in \citet{AbiJaberDeCarvalhoPham2024}. The present approach trades that generality for explicitness: restricted to the stationary regime, it expresses the adapted inverse of the friction directly through the factorization, as a closed-form filter. The two describe the same adapted optimum; Section~\ref{sec:gkgeneral} makes the identification, the iterative solution of the integral equation appearing as the Neumann expansion that the Gohberg--Krein factors sum in closed form.

% =====================================================================
\begin{table}[t]
\centering
\small
\begin{tabularx}{\textwidth}{@{}lX@{}}
\toprule
Symbol & Meaning \\
\midrule
\multicolumn{2}{@{}l}{\emph{Operators} (act on time-domain processes; no argument)}\\
$N,\ N_\pm$ & friction operator (position-referred) and its causal/anticausal Szeg\H{o} factors \\
$Q,\ Q_\pm$ & friction referred to the rate, $Q=N(-\partial_t^2)^{-1}$, and its factors (\S\ref{sec:fractional}) \\
$P_+$ & optional projection onto adapted processes \\
$I^\nu_\pm,\ D^\nu_\pm$ & Liouville fractional integrals and Marchaud fractional derivatives \\
$C_\pm$ & Gohberg--Krein (finite-horizon Volterra) factors \\
$G_T,\ I$ & rate-referred friction on $[0,T]$; identity \\
\addlinespace
\multicolumn{2}{@{}l}{\emph{Kernels and frequency-domain symbols} ($\hat{\cdot}$ denotes the Fourier transform)}\\
$g$ & transient-impact propagator kernel \\
$\hat n$ & friction symbol $\eta\omega^2+\gamma\hat g\,\omega^2+\lambda$; outer factors $\hat n_\pm$ \\
$\hat q=\hat n/\omega^2$ & rate-referred symbol $\eta+\gamma\hat g+\lambda/\omega^2$, with $\hat q_\pm=\hat n_\pm/(\mp i\omega)$ \\
$\hat\psi,\ \hat\varphi$ & outer spectral factors of the return $\mu$ ($S_\mu=|\hat\psi|^2$) and of the appreciation $\alpha$ (\S\ref{sec:fractional}) \\
$\pluspart{\,\cdot\,}$ & causal (plus) part of a symbol \\
\addlinespace
\multicolumn{2}{@{}l}{\emph{Filters and processes}}\\
$\hat\pi,\ \hat\chi=(-i\omega)\hat\pi$ & position filter and trade (rate) filter: innovations $\dot W\mapsto x^\star,\,u^\star$ \\
$\mu,\ \alpha$ & return signal (expected return) and expected remaining appreciation $\alpha_t=\E_t\int_t^\infty\mu_s\,ds$ \\
$u,\ x$ & trading rate and position, $x_t=\int_{-\infty}^t u_s\,ds$; forecast curve $\bar\mu(t,s)=\E_t[\mu_s]$ \\
$W,\ \dot W$ & Brownian motion and its innovations (white noise) \\
$\zeta,\ \xi^\star$ & whitened forecast $N_-^{-1}\bar\mu$; nonanticipativity multiplier \\
\addlinespace
\multicolumn{2}{@{}l}{\emph{Scalars}}\\
$\eta,\ \gamma,\ \lambda$ & temporary-impact weight, transient-impact weight, risk aversion \\
$\theta,\ \kappa$ & signal mean-reversion rate; impact resilience rate \\
$\beta,\ \nu$ & propagator exponent; fractional order $\nu=(1-\beta)/2$ \\
$c_\beta$ & power-law transform constant, $\hat g=c_\beta|\omega|^{\beta-1}$, $c_\beta=2\Gamma(1-\beta)\sin(\pi\beta/2)$ \\
$\Phi(\theta)$ & outer factor on the imaginary axis, $\hat n_+(i\theta)$ \\
$v,\ v_{\rm ant}$ & adapted value rate; anticipative (perfect-foresight) value rate \\
\bottomrule
\end{tabularx}
\caption{Notation. Operators are italic capitals and carry no argument; kernels are lowercase Latin; a hat denotes the Fourier transform; filters are lowercase Greek. The development is in the position variable; the rate variable is used for the scale-free kernel of Section~\ref{sec:fractional}, where the position is non-stationary, and on the finite horizon of Section~\ref{sec:boundary}.}
\label{tab:notation}
\end{table}

% =====================================================================
\section{The Wiener--Hopf solution}\label{sec:interior}

The stationary problem is solved in the three steps announced in \S\ref{sec:thesis}: factor the friction (\S\ref{sec:setup}), predict the signal (\S\ref{sec:meanrev}), and combine the two into a filter (\S\ref{sec:filter}). Throughout, fix a filtered probability space satisfying the usual conditions; Table~\ref{tab:notation} collects the notation. The return signal $\mu$ is adapted and mean-zero, with forecast curve
\[
\bar\mu(t,s) \;=\; \begin{cases}\mu_s, & s\le t,\\[2pt] \E_t[\mu_s], & s>t,\end{cases}
\]
decaying as $s\to\infty$, so that forecasts vanish at long horizons. The Fourier convention is $\hat f(\omega)=\int e^{i\omega t}f(t)\,dt$, under which causal kernels have transforms analytic in the upper half-plane. Admissible positions are adapted processes of finite friction energy $\E\langle x,Nx\rangle<\infty$; for $\eta=0$ this class admits positions with a Brownian component, whose trading rate carries white noise.

\subsection{Factorization of the friction}\label{sec:setup}

The symbol $\hat n$ of \eqref{eq:symbol} is positive and even, grows polynomially, and satisfies the Szeg\H{o} condition $\int|\log\hat n(\omega)|/(1+\omega^2)\,d\omega<\infty$, met by every kernel considered here, the pure power law included: its zero $\hat n\sim\gamma c_\beta|\omega|^{1+\beta}$ at $\omega=0$ is only logarithmic in $\log\hat n$, hence integrable \citep{Krein1962}. This log-integrability lets $N$ factor as $N=N_-N_+$ with symbol
\begin{equation}\label{eq:wh}
\hat n(\omega) = \hat n_-(\omega)\,\hat n_+(\omega), \qquad \hat n_- = \overline{\hat n_+},
\end{equation}
where $\hat n_+$ is analytic and zero-free in the upper half-plane. Then $N_+$ is causal with causal inverse, and $N_-=N_+^*$ is its anticausal adjoint. Conjugation by such a factor whitens the friction: under it the inner product $\langle\,\cdot\,,N\,\cdot\,\rangle$ becomes the flat $L^2$ product. The rate-referred factors carry one time-integration on each side, $\hat q_\pm=\hat n_\pm/(\mp i\omega)$. The outer factor $\hat n_+$ is given by the Szeg\H{o} formula
\begin{equation}\label{eq:szego}
\hat n_+(\omega) \;=\; \exp\!\Bigl(\frac{1}{2\pi i}\int_{-\infty}^{\infty}\Bigl(\frac{1}{t-\omega}-\frac{t}{1+t^2}\Bigr)\log \hat n(t)\,dt\Bigr), \qquad \operatorname{Im}\omega>0,
\end{equation}
unique up to a unimodular constant \citep{Krein1962}. Because $\hat n$ is even, evaluating \eqref{eq:szego} on the imaginary axis reduces to the real Poisson integral $\Phi(\theta):=\hat n_+(i\theta)$ of \eqref{eq:phi}, used repeatedly below.

\begin{assumption}[Friction]\label{ass:friction}
The weights satisfy $\eta,\gamma,\lambda\ge0$; the propagator $g$ is locally integrable and positive-definite, $\hat g\ge0$ in \eqref{eq:symbol}; and the symbol $\hat n$ is positive almost everywhere and satisfies the Szeg\H{o} condition, so that the factorization \eqref{eq:wh} exists.
\end{assumption}

\subsection{Prediction of the signal}\label{sec:meanrev}

The adapted trade responds to the signal's expected future. The trader cannot see the future path of $\mu$ and replaces it by the forecast curve $\bar\mu(t,s)=\E_t[\mu_s]$, $s\ge t$. For a Gaussian signal the forecast is a linear functional of the past; it reduces to a linear operator on the signal itself, or on a finite Markov lift, when the signal is finite-dimensional.

\begin{definition}[Markov signal]\label{def:mr}
A stationary signal is Markov of order $d$ if its forecast curve is a deterministic combination of the current value and its first $d-1$ derivatives,
\begin{equation}\label{eq:collapse}
\bar\mu(t,s) \;=\; \sum_{k=0}^{d-1}\mu^{(k)}_t\,g_k(s-t)\qquad (s\ge t),\qquad g_k(0)=\delta_{k0},
\end{equation}
equivalently if its spectral density is rational of degree $d$.
\end{definition}

Such a signal is a linear readout $\mu_t=c^{\!\top}X_t$ of a $d$-dimensional Markov state --- its lift; in observable coordinates the lift is $(\mu_t,\dot\mu_t,\dots,\mu^{(d-1)}_t)$, and \eqref{eq:collapse} carries the whole future in finitely many numbers. Consistency pins the profiles down: the tower property $\E_t[\bar\mu(t',s)]=\bar\mu(t,s)$ forces the mean dynamics to be linear, so the $g_k$ solve a $d$th-order constant-coefficient ODE --- exponential modes $e^{-r_j(s-t)}$, or damped oscillations when the rates $r_j$ are complex.

The mean-reverting case is $d=1$: the forecast is the current value times one profile, $\bar\mu(t,s)=\mu_t f(s-t)$, and the tower property forces $f(\tau)=e^{-\theta\tau}$ with $\theta>0$ the mean-reversion rate. This is the Ornstein--Uhlenbeck signal --- the only stationary Gaussian signal with a fixed forecast profile, by Doob's theorem --- and it also covers the non-Markov, non-Gaussian moving averages $\mu_t=\int_{-\infty}^t e^{-\theta(t-r)}\sigma_r\,dW_r$ with adapted stochastic volatility $\sigma$. For $d=2$ the forecast combines the signal and its velocity, a damped oscillation when the two rates are complex.

A signal with no finite-dimensional state has no such reduction: its forecast is the Wiener--Kolmogorov predictor, a linear operator on the whole past. Take $\mu$ stationary, Gaussian, purely non-deterministic, and generating its own filtration $\F_t=\sigma(\mu_s:s\le t)$. Its spectral density factors as $S_\mu=|\hat\psi|^2$ with $\hat\psi$ the outer (causal, zero-free) spectral factor, so $\mu=\psi\ast\dot W$ is a causal moving average of its innovations $\dot W$, the unit white noise generating the filtration --- the Wold representation \citep{Wiener1949,Whittle1963}. The forecast keeps the innovations up to $t$ and drops the later ones.

\subsection{The trading filter}\label{sec:filter}

The factorization and the forecast combine through the adapted inverse of $N$. The first-order condition \eqref{eq:foc} inverts the projected operator $P_+NP_+$, which the one-sided factors assemble. The lemma is stated for a general factored positive operator; it is applied to $N$ here, to the rate-referred $Q$ in Section~\ref{sec:fractional}, and to the finite-horizon operator in Section~\ref{sec:boundary}.

\begin{lemma}[Adapted projected inverse]\label{lem:pi}
Let $A=A_-A_+$ be a positive operator with $A_+$ causal with causal inverse and $A_-=A_+^*$ anticausal with anticausal inverse. Then, on adapted processes,
\begin{equation}\label{eq:pi}
(P_+\,A\,P_+)^{-1} \;=\; A_+^{-1}\,P_+\,A_-^{-1}.
\end{equation}
\end{lemma}

The proof, in Appendix~\ref{app:A}, is triangular bookkeeping: causality makes $A_+$ and its inverse preserve the adapted subspace, adjointness makes $A_-$ and its inverse preserve the orthogonal complement, and the projection annihilates what lands there. Among all square roots of $\hat n$, the factorization \eqref{eq:wh} selects the one whose whitening $N_-^{-1}$ leaves the space of processes invisible to adapted strategies invariant; this is the sense in which it is Cholesky with respect to the filtration.

\begin{remark}[Relation to operator factorization]\label{rmk:nest}
Lemma~\ref{lem:pi} is the projected-inverse identity behind the triangular factorization of a positive operator relative to a nest of subspaces \citep{GohbergKrein1970,Arveson1975}. Here the nest is the increasing family of adapted subspaces indexed by the filtration, and \eqref{eq:wh} is the corresponding filtration-indexed Cholesky factorization: the scalar Szeg\H{o} outer factor in the stationary case, and the Gohberg--Krein Volterra factor on a finite horizon (Section~\ref{sec:boundary}).
\end{remark}

\begin{theorem}[Optimal policy, general adapted signal]\label{thm:general}
Under Assumption~\ref{ass:friction}, the unique adapted optimum of \eqref{eq:objective} is
\begin{equation}\label{eq:policy}
x^\star \;=\; N_+^{-1}\,P_+\,N_-^{-1}\,\mu,
\qquad\text{equivalently}\qquad
x^\star_t = \bigl(N_+^{-1}\zeta\bigr)(t), \quad
\zeta_s = \bigl(N_-^{-1}\,\bar\mu(s,\cdot)\bigr)(s),
\end{equation}
where the second form uses the commutation of the deterministic operator $N_-^{-1}$ with conditional expectation on its argument.
\end{theorem}

The three factors in \eqref{eq:policy} are the three steps of the classical recipe: an anticausal whitening, the adapted projection, and a causal recolouring. The whitening would call for the future of the signal; the projection supplies instead its forecast (\S\ref{sec:meanrev}). Only the friction's factor enters the whitening; the signal's own spectral factor enters through the projection alone. Completing the square gives a second reading. Up to a constant the objective \eqref{eq:objective} equals $-\tfrac12\|x-N^{-1}\mu\|_N^2$ in the friction norm $\|y\|_N^2=\E\langle y,Ny\rangle$, so that $x^\star$ is the orthogonal projection, in that norm, of the anticipative optimum $N^{-1}\mu$ onto the adapted positions --- equivalently the causal Wiener--Kolmogorov estimate of the whitened return from the past.

\begin{remark}[Relation to Wiener--Kolmogorov prediction]\label{rmk:wk}
The construction is the Wiener--Hopf solution of a Wiener--Kolmogorov estimation problem, with the friction in place of the signal covariance. Linear prediction whitens a stationary signal by its own spectral factor $\hat\psi$, keeps the causal part, and recolours, returning the least-squares forecast from the past \citep{Wiener1949,Whittle1963}. The same three steps, carried out with the friction's factor $\hat n_+$ in place of $\hat\psi$, give the optimal position: by the estimation reading above, $x^\star$ is the causal estimate of the anticipative optimum $N^{-1}\mu$ in the friction metric. Prediction and trading are thus one construction on two quadratic forms, the signal covariance and the friction --- the correspondence that carries frequency-domain factorization from filtering into linear--quadratic control \citep{Whiteman1985}.
\end{remark}

\begin{remark}[The nonanticipativity multiplier]\label{rmk:duality}
Substituting \eqref{eq:policy} into $Nx^\star$ writes the first-order condition \eqref{eq:foc} without the projection,
\begin{equation}\label{eq:multiplier}
N\,x^\star = \mu - \xi^\star, \qquad \xi^\star = N_-\,(I-P_+)\,N_-^{-1}\mu, \qquad P_+\xi^\star = 0,
\end{equation}
so the adapted optimum solves the unconstrained first-order condition for the modified return $\mu-\xi^\star$. A multiplier with vanishing optional projection is the defining property of the nonanticipativity multipliers of \citet{RockafellarWets1976}, established there for general convex multistage problems; \eqref{eq:multiplier} is their closed form here. Because $\E\langle x,\xi^\star\rangle=0$ for adapted $x$, the value forgone to adaptedness equals the friction energy $\tfrac12\E\langle\xi^\star,N^{-1}\xi^\star\rangle$ of the multiplier, per unit time in the stationary case --- the shadow price of information of \S\ref{sec:value} below.
\end{remark}

\begin{assumption}[Signal]\label{ass:signal}
The return signal $\mu$ is stationary, Gaussian, purely non-deterministic, and generates its own filtration, with Wold representation $\mu=\psi\ast\dot W$ and spectral density $S_\mu=|\hat\psi|^2$, $\hat\psi$ the outer spectral factor.
\end{assumption}

For a stationary signal the projection is itself a filter. In the innovations $\dot W$ the whitened return has symbol $h:=\hat\psi\,\hat n_-^{-1}$; since conditioning at time $s$ discards the innovations after $s$, the projection truncates the moving-average kernel at lag zero and acts on symbols as the Riesz plus-part $\pluspart{\,\cdot\,}$ onto non-negative lags. Then \eqref{eq:policy} is a transfer function.

\begin{theorem}[Optimal trading filter]\label{thm:filter}
Under Assumptions~\ref{ass:friction} and~\ref{ass:signal}, and $\lambda>0$, the optimal position is the linear filter $x^\star=\pi\ast\dot W$ with transfer function
\begin{equation}\label{eq:filter}
\hat\pi \;=\; \hat n_+^{-1}\,\pluspart{\,h\,},
\qquad h = \hat\psi\,\hat n_-^{-1},
\end{equation}
the optimal rate the filter $\hat\chi=(-i\omega)\,\hat\pi$, and the attained value rate, the expected profit per unit time, is
\begin{equation}\label{eq:value}
v \;=\; \frac{1}{4\pi}\,\bigl\|\pluspart{h}\bigr\|_{L^2}^2 .
\end{equation}
\end{theorem}

For a Markov signal the same filter reads on the observed signal rather than the innovations. The whitened forecast then reduces to a differential operator: with $\rho_k=(N_-^{-1}g_k)(0)$,
\begin{equation}\label{eq:zeta-collapse}
\zeta_s = (N_-^{-1}\bar\mu(s,\cdot))(s) = \sum_{k=0}^{d-1}\rho_k\,\mu^{(k)}_s = P(\partial)\,\mu_s,
\end{equation}
so that the policy filters the signal and its derivatives, $\hat x^\star(\omega)=P(-i\omega)\,\hat\mu(\omega)/\hat n_+(\omega)$, a rational filter. For the Ornstein--Uhlenbeck signal ($d=1$) the single pole at $\omega=-i\theta$ gives $P(\partial)=\rho=1/\hat n_-(-i\theta)=1/\Phi(\theta)$ (Appendix~\ref{app:B}), with
\begin{equation}\label{eq:phi}
\Phi(\theta) \;:=\; \hat n_+(i\theta) \;=\; \exp\Bigl[\frac{\theta}{2\pi}\int\frac{\log\hat n(t)}{\theta^2+t^2}\,dt\Bigr] \;>\;0,
\end{equation}
the outer factor of $\hat n$ on the imaginary axis, so the position filter is the single formula covering all three frictions,
\begin{equation}\label{eq:ou-filter}
\hat x^\star(\omega) \;=\; H(\omega)\,\hat\mu(\omega), \qquad H(\omega)=\frac{1}{\Phi(\theta)\,\hat n_+(\omega)},
\end{equation}
with value $v=\sigma^2\theta/4\Phi(\theta)^2$, where $\theta^2\sigma^2$ is the innovation variance of $\mu$ ($\sigma^2$ that of $\alpha=\mu/\theta$). For $d=2$ the filter uses the signal and its velocity.

The hypothesis $\lambda>0$ keeps the whitened return and the position filter in $L^2$: without risk aversion $\hat n$ vanishes at $\omega=0$ and the position is non-stationary --- the regime of Section~\ref{sec:fractional}, where the theory is referred to the rate. The rate filter $\hat\chi$ lies in $L^2$ under the spectral-decay hypothesis of Appendix~\ref{app:B}, and otherwise, notably for the exponential kernel without temporary cost, carries a lag-zero white-noise atom.

The value \eqref{eq:value} has a benchmark. The perfectly-anticipating trader replaces the projection $\pluspart{\,\cdot\,}$ by the identity, extracting
\begin{equation}\label{eq:vant}
v_{\rm ant} \;=\; \frac{1}{4\pi}\int \frac{S_\mu}{\hat n}\,d\omega,
\end{equation}
the value attainable with the whole signal path in hand, finite for $\lambda>0$. The \emph{causality gap} $v/v_{\rm ant}\le1$ is the fraction of that clairvoyant value the adapted trader keeps; the shortfall $v_{\rm ant}-v$ equals, by \eqref{eq:value}, the $L^2$ energy of the anticausal part of the half-whitened signal $h=\hat\psi\hat n_-^{-1}$, and by Remark~\ref{rmk:duality} the friction energy of the nonanticipativity multiplier $\xi^\star$ --- the shadow price of adaptedness, in closed form. Under power-law impact the gap takes a scale-free value, $\sin(\pi\beta/2)$ (\S\ref{sec:value}).

% =====================================================================
\section{The finite-horizon factorization}\label{sec:boundary}

The stationary theory describes an ongoing operation on the whole line. Most execution problems are posed instead on a fixed window $[0,T]$: a broker works a parent order over a session, or liquidates an inherited position by a deadline, ending flat, $x_T=0$. Two features are then absent from the stationary theory: the two endpoints of the window, and the terminal inventory constraint. On $[0,T]$ translation invariance is lost, the signal and friction no longer share a Fourier basis, and the Szeg\H{o} formula does not apply.

The factorization survives; only its computation changes. On the window the working variable is again the rate, which carries the terminal constraint and the endpoint blocks; the rate-referred friction on $[0,T]$, $G_T=\eta I$ plus a continuous-kernel integral operator, is positive. Relative to the nest of adapted subspaces such an operator factors triangularly, $G_T=C_-C_+$ with $C_+$ a causal Volterra operator and $C_-=C_+^*$ \citep{GohbergKrein1970,Arveson1975}. Lemma~\ref{lem:pi} applies verbatim, and the optimum is
\begin{equation}\label{eq:finiteT}
u^\star \;=\; C_+^{-1}\,P_+\,C_-^{-1}\,\alpha^{\rm eff},
\qquad \alpha^{\rm eff} = \alpha + \sum_k \xi_k e_k,
\end{equation}
where $\alpha^{\rm eff}$ adjoins a multiplier $\xi_k$ for each linear position constraint $\langle e_k,x\rangle=0$, with a process-valued multiplier for the pathwise liquidation constraint $x_T=0$; each modifies the signal exactly as the adaptedness constraint does in \eqref{eq:multiplier}, confined to its constraint's annihilator, as in the constrained propagator problem of \citet{AbiJaberDeCarvalhoPham2024}. The Szeg\H{o} formula is replaced by a Volterra kernel, and nothing else in the construction moves.

\subsection{Relaxation to the stationary filter}\label{sec:bdrylayer}

The two endpoints and the terminal constraint act locally: their influence decays into the interior, where the finite-horizon policy relaxes to the stationary filter of \S\ref{sec:filter}. The bulk of a long session is governed by the whole-line solution, and only the start-up and terminal windows need separate treatment.

\begin{proposition}[Boundary-layer decay]\label{prop:boundary}
Let $u^{\star,T}$ be the finite-horizon optimum on $[0,T]$, $u^\star$ the whole-line stationary optimum for the same bounded signal, and $d(t)=\min(t,T-t)$ the distance to the nearest endpoint. Then, for constants $C(\beta)$, $C$ depending only on the kernel parameters,
\begin{equation}\label{eq:bdry-decay}
\bigl|u^{\star,T}_t - u^\star_t\bigr| \;\le\;
\begin{cases}
C(\beta)\,\|\alpha\|_\infty\,d(t)^{-\nu}, & \text{power-law kernel},\ \nu=\tfrac{1-\beta}{2},\\[3pt]
C\,\|\alpha\|_\infty\,e^{-b_1 d(t)}, & \text{rational symbol, } b_1 \text{ the slowest zero of } \hat n_+ .
\end{cases}
\end{equation}
\end{proposition}

The proof is in Appendix~\ref{app:C}. The boundary layer has width $\sim1/b_1$ for a rational kernel and is set by the impact memory for a power law; it is thin except for the small $\nu$ of the empirical power-law exponents, a quantitative statement of how long-memory impact makes horizons matter. Figure~\ref{fig:bdry} shows the finite-horizon rate tracking the stationary filter through the interior.

\begin{figure}[t]
\centering
\includegraphics[width=0.72\linewidth]{fig_boundary_layer}
\caption{Finite horizon versus whole line ($\eta=0.5$, $\gamma=1$, $\kappa=2$, $\lambda=1$, $T=20$). The finite-horizon optimal rate tracks the whole-line stationary rate through the interior and departs only in the start-up and terminal boundary layers, as quantified by Proposition~\ref{prop:boundary}; the shading marks the boundary-layer scale, not a fitted width.}
\label{fig:bdry}
\end{figure}

\subsection{General kernels}\label{sec:gkgeneral}

The explicit Volterra factor depends on the kernel, and closes in two classes. For a rational symbol it is a finite sum of endpoint-anchored exponential modes matched at the two boundaries, giving the block-plus-continuous liquidation of \S\ref{sec:expFH}; for a scale-free kernel it is a conjugated fractional factor, giving the U-shaped and Fredholm profiles of \S\ref{sec:powerlawFH}. Between the two the factor kernel solves the Gohberg--Krein integral equations on $[0,T]$ \citep{GohbergKrein1970,Arveson1975} and is computed numerically.

The factorization nonetheless represents the general-propagator solution in closed operator form. The adapted optimum of \eqref{eq:objective} is unique (Theorem~\ref{thm:general}), so the Gohberg--Krein policy \eqref{eq:finiteT} and the general-propagator solutions of the literature are one process written several ways: the free-boundary backward SDE and operator Riccati equation of \citet{AbiJaberNeuman2025}, the operator resolvents of \citet{AbiJaberNeumanTuschmann2024}, and the stochastic Fredholm equation with constraint multipliers of \citet{AbiJaberDeCarvalhoPham2024}. On $[0,T]$ the identification is term by term: the iterative solution of the stochastic Fredholm equation is the Neumann expansion that the Gohberg--Krein factors sum in closed form. In the stationary regime the Szeg\H{o} factor gives the closed-form filter of \S\ref{sec:filter}, and in the interior the finite-horizon solution relaxes to it (Proposition~\ref{prop:boundary}).

% =====================================================================
\section{Exponential-kernel solutions}\label{sec:recovery}

The exponential and rational frictions --- inventory risk, temporary cost, exponential resilience, and their combinations --- have a rational symbol, so the Szeg\H{o} factor is a ratio of polynomials and the optimal position is a finite sum of exponential moving averages of the return signal: their rates are the zeros of $\hat n_+$, their number the number of poles, and risk aversion sets their scale. The classical portfolio and execution rules are the low-order cases, on the whole line (\S\ref{sec:rational}) and on a finite horizon (\S\ref{sec:expFH}).

\subsection{Stationary filters}\label{sec:rational}

\emph{Risk alone: the Markowitz rule.} With $\eta=\gamma=0$ the symbol is the constant $\hat n\equiv\lambda$, the factorization is trivial, and the \emph{position} is the memoryless mean--variance rule $x^\star_t=\mu_t/\lambda=\E_t[-\dot\alpha_t]/\lambda$ of \citet{Markowitz1952,Merton1971}, the expected return over risk aversion. The position responds to the return $\mu$, whereas under temporary impact alone (Section~\ref{sec:fractional}) the trading \emph{rate} is proportional to the appreciation signal $\alpha$. For an Ornstein--Uhlenbeck signal $\mu_t=\theta\alpha_t$, and the value is $v=\theta\sigma^2/4\lambda$.

\emph{Temporary cost with risk: the aim portfolio.} With $\gamma=0$ the friction is the differential operator $N=\lambda-\eta\,\partial_t^2$, local like the Markowitz case, $\hat n_+=\sqrt\eta\,(a-i\omega)$ with $a=\sqrt{\lambda/\eta}$: one moving average, and differentiation puts the filter in partial-adjustment form,
\begin{equation}\label{eq:gp}
u^\star_t = a\,(\mathrm{aim}_t - x^\star_t), \qquad
\mathrm{aim}_t = \frac{a}{a+\theta}\cdot\frac{\theta\alpha_t}{\lambda}:
\end{equation}
trade at speed $a$ toward an aim that discounts the Markowitz position by the persistence factor $a/(a+\theta)$. This is the continuous-time aim portfolio of \citet{GarleanuPedersen2016} (discrete time, \citet{GarleanuPedersen2013}; general targets, \citet{BankSonerVoss2017}): a diagonal friction gives exactly one pole.

\emph{Exponential propagator with risk.} For $g=e^{-\kappa|t|}$ and $\eta=0$, $\hat n_+=\sqrt A\,(m-i\omega)/(\kappa-i\omega)$ with $A=2\kappa\gamma+\lambda$, $m=\kappa\sqrt{\lambda/A}$, and \eqref{eq:ou-filter} becomes the current signal plus one moving average,
\begin{equation}\label{eq:ema}
x^\star_t \;=\; \frac{\kappa+\theta}{A(m+\theta)}\Bigl[\mu_t + (\kappa-m)\!\int_{-\infty}^t e^{-m(t-s)}\mu_s\,ds\Bigr],
\end{equation}
whose rate $m$ travels from $0$ to $\kappa$ as risk aversion grows. Exponential resilience with an Ornstein--Uhlenbeck signal is the model of \citet{LehalleNeuman2019}, whose explicit solution is the risk-free, singular finite-horizon liquidation generalizing \citet{ObizhaevaWang2013}; \eqref{eq:ema} is its stationary, no-temporary-cost form, risk aversion holding the position in place of the horizon.

\emph{Exponential propagator with temporary cost and risk.} With all three frictions the symbol is rational, and the outer factor $\hat n_+=\sqrt\eta\,(b_1-i\omega)(b_2-i\omega)/(\kappa-i\omega)$ has two zeros ($b_1^2b_2^2=\lambda\kappa^2/\eta$, $b_1^2+b_2^2=\kappa^2+(2\kappa\gamma+\lambda)/\eta$), so \eqref{eq:ou-filter} becomes a two-average filter, $x^\star_t=\frac{1}{\Phi(\theta)\sqrt\eta}\sum_i w_i\!\int_{-\infty}^t e^{-b_i(t-s)}\mu_s\,ds$ with $w_i=(\kappa-b_i)/(b_j-b_i)$. \citet{NeumanVoss2022} solve this by a system of coupled forward--backward SDEs and give the strategy as a feedback control; its gains reduce, in the stationary regime, to these two rates and two weights --- the zeros and residues of the symbol --- and the finite-horizon gains converge to them in the interior, the closed-loop poles equalling the moving-average rates $b_1,b_2$. Both degenerations are exact: at $\gamma=0$ the rates are $\{a,\kappa\}$ and the $\kappa$-pole cancels the kernel zero, recovering \eqref{eq:gp}; as $\eta\to0$, $b_1\to m$ and $\sqrt\eta\,(b_2-i\omega)\to\sqrt A$, recovering \eqref{eq:ema}. A rational friction with $n$ poles gives a filter with at most $n$ averages. Figure~\ref{fig:structure} shows how these rates and the resulting policy shape depend on the friction.

\begin{figure}[t]
\centering
\includegraphics[width=\linewidth]{fig_filter_structure}
\caption{The optimal policy across kernels and parameters. (a)~For the temporary-plus-resilient kernel ($\eta=0.5$, $\gamma=1$, $\kappa=2$), the two moving-average rates $b_1,b_2$ against risk aversion: the slow rate $b_1\to0$ as $\lambda\to0$ (the position becomes non-stationary without risk) and both grow with $\lambda$. (b)~The position impulse response $g_x(\tau)$ by kernel ($\theta=1$, normalized at $\tau=0.3$): the exponential and two-average kernels have finite memory (exponential cutoff), the power-law kernel a heavy tail (Section~\ref{sec:fractional}), and the Markowitz limit is an instantaneous spike. (c)~The power-law policy across $\beta$: a scale-free kernel of slope $(\beta-1)/2$.}
\label{fig:structure}
\end{figure}

\subsection{Finite-horizon liquidation: block-plus-continuous}\label{sec:expFH}

On a finite horizon the exponential kernel places a discrete block at each endpoint and trades continuously between --- the block-plus-continuous liquidation of \citet{ObizhaevaWang2013}, the rational finite-horizon counterpart of \eqref{eq:ema}. Its two blocks reflect the finite impact at zero lag, $g(0)<\infty$: instantaneous trading costs only finitely, so the optimum pays for it at the two boundaries. The finite-horizon factor is the finite sum of endpoint-anchored exponential modes matched at the two boundaries (\S\ref{sec:gkgeneral}); risk aversion holds the position where the horizon does in the risk-free liquidation. The power law, with $g(0)=\infty$, forbids the blocks, and its liquidation is continuous (\S\ref{sec:powerlawFH}).

% =====================================================================
\section{The power-law kernel}\label{sec:fractional}

The power-law kernel is scale-free: its transform $\hat g(\omega)=c_\beta|\omega|^{\beta-1}$ diverges as $\omega\to0$, so the impact is not integrable and has no characteristic timescale. It is a superposition of exponentials over all decay rates, $|t|^{-\beta}=\Gamma(\beta)^{-1}\int_0^\infty e^{-r|t|}r^{\beta-1}\,dr$: a continuum where a rational kernel has finitely many, and the reason the optimal filter is fractional. This is long memory, in contrast with the exponential kernel $e^{-\kappa|t|}$, whose impact is integrable and dies on the single scale $1/\kappa$. The optimum reflects this memory in two ways: the policy is a fractional derivative (\S\ref{sec:fracpolicy}), and the value of anticipation obeys the scale-free law $\sin(\pi\beta/2)$ (\S\ref{sec:value}).

With $\eta=\lambda=0$ the friction symbol $\hat n=\gamma c_\beta|\omega|^{1+\beta}$ vanishes at $\omega=0$: the optimal position is non-stationary, and the stationary object is the rate. This section therefore works in the rate variable and the appreciation signal $\alpha_t=\E_t\int_t^\infty\mu_s\,ds$ of \S\ref{sec:problem}, with forecast curve $\bar\alpha(t,s)=\E_t[\alpha_s]=\int_s^\infty\bar\mu(t,r)\,dr$; for the Ornstein--Uhlenbeck signal $\alpha=\mu/\theta$. Referred to the rate the friction is $Q$ with symbol $\hat q=\hat n/\omega^2=\gamma\hat g$ and factors $\hat q_\pm=\hat n_\pm/(\mp i\omega)$; Lemma~\ref{lem:pi} applies to $Q=Q_-Q_+$ as it does to $N$, and the policy of Theorem~\ref{thm:general} reads
\begin{equation}\label{eq:rate-policy}
u^\star \;=\; Q_+^{-1}\,P_+\,Q_-^{-1}\,\alpha,
\qquad
u^\star_t=\bigl(Q_+^{-1}\zeta\bigr)(t),\quad \zeta_s=\bigl(Q_-^{-1}\bar\alpha(s,\cdot)\bigr)(s),
\end{equation}
the same three operations, with the whitened forecasts equal, $\bigl(Q_-^{-1}\bar\alpha(s,\cdot)\bigr)(s)=\bigl(N_-^{-1}\bar\mu(s,\cdot)\bigr)(s)$ (Appendix~\ref{app:A}). For a stationary signal the filter and value of Theorem~\ref{thm:filter} likewise refer to the rate: $\hat\chi=\hat q_+^{-1}\pluspart{h_\alpha}$ and $v=\frac{1}{4\pi}\|\pluspart{h_\alpha}\|^2$ with $h_\alpha:=\hat q_-^{-1}\hat\varphi$, $\hat\varphi$ the outer spectral factor of $\alpha$ ($S_\alpha=|\hat\varphi|^2$); these remain in $L^2$ at $\lambda=0$ (Appendix~\ref{app:B}).

\subsection{The policy is a fractional derivative}\label{sec:fracpolicy}

At one extreme, temporary cost alone ($\gamma=\lambda=0$) makes the friction local: the optimal rate is the memoryless myopic rule $u^\star_t=\alpha_t/\eta$, proportional to the appreciation signal $\alpha$ --- whereas Markowitz makes the \emph{position} proportional to the return $\mu=\E_t[-\dot\alpha]$ (Section~\ref{sec:recovery}). Power-law impact is the opposite extreme: its policy is a fractional derivative of the signal. The fractional integral of order $\nu\in(0,1)$ on the whole line is the Liouville operator
\[
(I^\nu_+f)(t)=\frac{1}{\Gamma(\nu)}\int_{-\infty}^t(t-s)^{\nu-1}f(s)\,ds,
\]
with $I^\nu_-$ its time reflection; it interpolates between the identity ($\nu\to0$) and a single integration ($\nu=1$), weighting the past of $f$ by the power-law kernel $(t-s)^{\nu-1}$. Its inverse $D^\nu_\pm=(I^\nu_\pm)^{-1}$, the fractional derivative of order $\nu$, is non-local, determined by the entire history rather than the local slope; for the bounded signals here it takes the Marchaud form \citep{SamkoKilbasMarichev1993},
\[
(D^\nu_\pm f)(t)=\frac{\nu}{\Gamma(1-\nu)}\int_0^\infty\frac{f(t)-f(t\mp s)}{s^{1+\nu}}\,ds,
\]
whose increment regularizes the weight $s^{-1-\nu}$ and keeps the integral finite for bounded, non-decaying $f$. The power-law symbol is a fractional power of frequency: for $g(t)=|t|^{-\beta}$ the transform is $\hat g(\omega)=c_\beta|\omega|^{\beta-1}$, so the Wiener--Hopf factors are fractional integrals, $Q_\pm=(\gamma c_\beta)^{1/2}I^\nu_\pm$ with $\nu=(1-\beta)/2$, and their inverses are the fractional derivatives $D^\nu_\pm$. The rate policy \eqref{eq:rate-policy} then reads
\begin{equation}\label{eq:fractional}
u^\star_t \;=\; \frac{1}{\gamma c_\beta}\,(D_+^\nu\zeta)(t), \qquad
\zeta_s = \bigl(D_-^\nu\,\bar\alpha(s,\cdot)\bigr)(s):
\end{equation}
an anticausal fractional derivative of the forecast curve, followed by a causal fractional derivative of the whitened forecast, both of order $\nu\in(0.2,0.4)$ for the empirical range of $\beta$. For an Ornstein--Uhlenbeck signal $\zeta_s=\theta^\nu\alpha_s$ and $u^\star=(\theta^\nu/\gamma c_\beta)D^\nu_+\alpha$, with value $v=\sigma^2\theta^{-\beta}/4\gamma c_\beta$; comparing with the anticipative value \eqref{eq:vant} gives the causality gap \eqref{eq:sinlaw}. The filter has long memory: the Marchaud representation weights past signal increments by the slowly decaying kernel $(t-s)^{-1-\nu}$, and on a grid $D^\nu_+$ is a Toeplitz operator whose Gr\"unwald--Letnikov weights inherit this tail.

The position that this rate accumulates is a fractional \emph{integral} of the return. The position filter \eqref{eq:ou-filter}, $\hat x^\star=\hat\mu/(\Phi(\theta)\hat n_+)$, has $\hat n_+\propto(-i\omega)^{(1+\beta)/2}$ for the pure power law, so the position is a causal fractional integral of the return of order $(1+\beta)/2$,
\begin{equation}\label{eq:fracpos}
\hat x^\star(\omega) \;\propto\; (-i\omega)^{-(1+\beta)/2}\,\hat\mu(\omega), \qquad x^\star \propto I^{(1+\beta)/2}_+\mu,
\end{equation}
the scale-free, long-memory counterpart of the memoryless Markowitz position $x^\star=\mu/\lambda$. Low frequencies are amplified: the squared gain is $|\omega|^{-(1+\beta)}$, so the position has spectral density $|\omega|^{-(1+\beta)}S_\mu(\omega)$ up to a constant. Since $1+\beta>1$, this diverges at the origin unless the return's low-frequency power is depleted: the position is stationary if and only if $\int|\omega|^{-(1+\beta)}S_\mu(\omega)\,d\omega<\infty$, that is, if $S_\mu(\omega)$ vanishes at $\omega=0$ faster than $|\omega|^\beta$. A persistent return fails the test: for the Ornstein--Uhlenbeck signal $S_\mu(0)\ne0$, and the position drifts without bound, a fractional analogue of the random walk that integrating a persistent return produces. This is why the pure power-law theory is referred to the rate, stationary for every signal, rather than to the position; a risk term $\lambda>0$ confines the position for any return (\S\ref{sec:filter}), at the cost of leaving the scale-free regime.

Between the two ends the position filter $H(\omega)=[\Phi(\theta)\hat n_+(\omega)]^{-1}$ has magnitude proportional to $\hat n(\omega)^{-1/2}$, interpolating across the crossover frequencies of $\hat q$: risk-dominated at low frequency, transient-impact-dominated in the middle, and temporary-cost-dominated at high frequency, the crossovers falling at $\omega_c=(\lambda/\gamma c_\beta)^{1/(1+\beta)}$ and $\omega_*=(\gamma c_\beta/\eta)^{1/(1-\beta)}$. Figure~\ref{fig:transfer} places the two transient kernels side by side in both domains: the exponential settles on a finite-memory plateau and an impulse response that cuts off, while the power law slides without a plateau and leaves an algebraic tail, the temporary cost capping the short lag and the risk term the long.

\begin{figure}[t]
\centering
\includegraphics[width=\linewidth]{fig_transfer_impulse}
\caption{Anatomy of the trading filter, exponential versus power-law transient impact ($\gamma=1$, $\kappa=2$, $\beta=0.5$; small $\lambda=0.1$, $\eta=0.05$). (a)~Transfer function $|H(\omega)|=1/\sqrt{\hat n(\omega)}$: both share the low-frequency Markowitz plateau $1/\sqrt\lambda$ (edge $\omega_c$) and roll off as $\omega^{-1}$ under the temporary cost at high frequency (edges $\omega_e$ for the exponential, $\omega_*$ for the power law); between them the finite-memory exponential settles on the plateau $1/\sqrt{2\kappa\gamma}$ while the long-memory power law slides as $\omega^{-(1+\beta)/2}$. (b)~Impulse response $h(\tau)=\mathcal{F}^{-1}[\hat n_+^{-1}]$: the temporary cost caps the short lag at $1/\sqrt\eta$; past it the exponential filter cuts off (short memory) and the power law leaves an algebraic tail (long memory).}
\label{fig:transfer}
\end{figure}

\subsection{The value of information}\label{sec:value}

Under pure power-law impact the causality gap of \S\ref{sec:filter} is scale-free, independent of the signal speed. For the Ornstein--Uhlenbeck signal the adapted value is $v=\sigma^2\theta^{-\beta}/4\gamma c_\beta$ (\S\ref{sec:fracpolicy}); the anticipative benchmark \eqref{eq:vant} is referred to the rate at $\lambda=0$, $v_{\rm ant}=\frac{1}{4\pi}\int S_\alpha/\hat q\,d\omega$, finite where the position form diverges (Appendix~\ref{app:B}), and their ratio is
\begin{equation}\label{eq:sinlaw}
\frac{v}{v_{\rm ant}} \;=\; \sin\!\frac{\pi\beta}{2}.
\end{equation}
Adaptedness is nearly free for short-memory impact ($\beta\to1$) and forfeits almost all the value for near-permanent impact ($\beta\to0$); at the empirical $\beta\in(0.2,0.6)$ the adapted trader keeps between a third and four-fifths of the clairvoyant value (Figure~\ref{fig:gap}).

\begin{figure}[t]
\centering
\includegraphics[width=0.62\linewidth]{fig_causality_gap}
\caption{The cost of causality. Under power-law impact the adapted value is a fixed fraction of the clairvoyant value, $v/v_{\rm ant}=\sin(\pi\beta/2)$ of~\eqref{eq:sinlaw}, independent of the signal speed: adaptedness is nearly free for short-memory impact ($\beta\to1$) and forfeits almost all the value for near-permanent impact ($\beta\to0$), keeping between a third and four-fifths of the clairvoyant value over the empirical range (shaded).}
\label{fig:gap}
\end{figure}

\subsection{Finite-horizon liquidation: the U-shape}\label{sec:powerlawFH}

On a finite horizon the scale-free kernel has an explicit Volterra factor. For the power-law kernel ($\eta=\lambda=0$) the causal factor $C_+$ is the terminal-anchored Volterra operator with kernel
\begin{equation}\label{eq:gk-kernel}
c_+(t,s) = (\gamma c_\beta)^{1/2}\,\Bigl(\frac{T-s}{T-t}\Bigr)^{\!\nu}\frac{(t-s)^{\nu-1}}{\Gamma(\nu)},
\qquad 0\le s\le t\le T,
\end{equation}
verified by direct kernel integration, $C_-C_+=G_T$ including the constant. This is the whole-line fractional factor $I^\nu_+$ conjugated by multiplication by $(T-t)^\nu$: the terminal weight anchors the convolution factor to the endpoint $T$, and homogeneity of the kernel is what allows the anchoring to be a scalar weight. Each factor sees one endpoint, $C_+$ the terminal and its time reflection $C_-$ the initial; the identity \eqref{eq:pi} places the causal factor on the right of \eqref{eq:finiteT}, so the adapted projection sits between the factors and forces $C_+$ to anchor at $T$. Away from the endpoints the weight tends to one and $C_+$ tends to the whole-line factor, which is Proposition~\ref{prop:boundary} for this kernel.

The liquidation profiles of the execution literature follow. For a constant signal the projection is trivial and \eqref{eq:finiteT} reduces to $G_T^{-1}\mathbf 1$ up to the terminal-inventory multiplier, the U-shaped profile $u^\star(t)\propto[t(T-t)]^{(\beta-1)/2}$ of \citet[Ex.~2.30]{GatheralSchiedSlynko2012} --- continuous, without the endpoint blocks of the exponential kernel (\S\ref{sec:expFH}), because $g(0)=\infty$. For a Gaussian Volterra signal, the first-order condition of the finite-horizon problem is the stochastic Fredholm equation that \citet{FordeSanchezSmith2022} solve on the Wiener chaos; the time reflection of \eqref{eq:gk-kernel} is the left-anchored outer factor they use to invert $G_T$ in full, where the factor order is immaterial. Formula \eqref{eq:finiteT} is its solution operator, and expanding it on the chaos reproduces their Fredholm systems, with the generalized Abel inversion of \citet{ChakrabartiGeorge1994} implementing $G_T^{-1}$ and weighted Riemann--Liouville operators the individual factors.

% =====================================================================
\section{Concluding remarks}\label{sec:conclusion}

Trading against impact is a problem of timing. Unconstrained, the first-order condition is a single inversion of the friction operator; under adaptedness, with a friction that couples different times, the trader cannot truncate that inversion, because projection onto the past does not commute with the inverse. Factoring the friction along the filtration, into a causal part and its anticausal adjoint, resolves this, and the adapted inverse assembles from the one-sided factors --- whiten the forecast, project, recolour.

On the whole line the factor is a frequency-domain formula and the policy is a \emph{filter}, a transfer function that maps the observed signal to the position; the classical portfolio and execution rules are its special cases --- Markowitz when the friction is constant, the aim portfolio when it is a second-order differential operator, the moving-average filters of the resilience models, and a fractional derivative when the impact is scale-free. The shape of the policy is set by the impact memory: a rational kernel has finitely many relaxation rates and the position is a finite sum of moving averages; a power-law kernel has a continuum and the policy is a fractional derivative, keeping a fraction $\sin(\pi\beta/2)$ of the value a clairvoyant trader would earn.

Three restrictions set the scope. The costs are quadratic in the rate and position, so nonlinear impact lies outside the model. The explicit filters require a stationary Gaussian signal; for a general adapted signal, Theorem~\ref{thm:general} gives the policy through the forecast curve but not as a transfer function. And the asset is one-dimensional: a portfolio version calls for a matrix Wiener--Hopf factorization, explicit only in special cases. On finite horizons the theory is complete up to the boundary terms attached to constraints, for which \eqref{eq:finiteT} gives the operator form but not yet sharp constants.

% =====================================================================
\appendix

\section{Proof of Lemma~\ref{lem:pi} and Theorem~\ref{thm:general}}\label{app:A}
Causality of $A_+$ gives $(I-P_+)A_+P_+=0$, so $A_+$ and its causal inverse preserve the adapted subspace; taking adjoints and using $A_-=A_+^*$ gives $P_+A_-(I-P_+)=0$, so $A_-$ and $A_-^{-1}$ preserve the orthogonal complement. For adapted $v$, write $A_-v=P_+A_-v+(I-P_+)A_-v$; applying $A_-^{-1}$, the image of the second term stays in the complement and is annihilated by $P_+$, so $P_+A_-^{-1}P_+A_-v=P_+v=v$. Hence for adapted $u$,
\[
A_+^{-1}P_+A_-^{-1}\,(P_+AP_+)\,u = A_+^{-1}\bigl[P_+A_-^{-1}P_+A_-\bigr]A_+u = A_+^{-1}A_+u = u,
\]
using $P_+A_+u=A_+u$. The reverse composition equals the identity by the same bookkeeping, or from positivity of $P_+AP_+$ on the adapted subspace, for which a one-sided inverse is two-sided. For the pure power-law kernel the factors are unbounded and the computation holds on the dense domain of signals with $\int(1+\omega^2)S_\alpha(\omega)/\hat q(\omega)\,d\omega<\infty$ (Appendix~\ref{app:B}); any $\eta>0$ or $\lambda>0$ gives bounded inverses.

Theorem~\ref{thm:general} follows with $A=N$: strict convexity ($\hat n>0$ almost everywhere) makes \eqref{eq:foc} necessary and sufficient with a unique solution, which Lemma~\ref{lem:pi} inverts. The forecast-curve form follows from conditional Fubini: $N_-^{-1}$ has a deterministic anticausal kernel, so $\E_s[(N_-^{-1}\mu)(s)]$ integrates that kernel against $\E_s[\mu_r]=\bar\mu(s,r)$ for $r\ge s$. The rate form \eqref{eq:rate-policy} follows by differentiating \eqref{eq:policy}: $\partial_t\circ N_+^{-1}=Q_+^{-1}$, and on forecast curves $Q_-^{-1}=N_-^{-1}\circ(-\partial_s)$ with $-\partial_s\bar\alpha(s,\cdot)=\bar\mu(s,\cdot)$, so the whitened forecasts coincide.

\section{Proof of Theorem~\ref{thm:filter} and the Ornstein--Uhlenbeck formulas}\label{app:B}
The whitened return $N_-^{-1}\mu$ is the stationary filter of $\dot W$ with symbol $h=\hat\psi\,\hat n_-^{-1}$. Conditional expectation at $s$ annihilates innovations after $s$, so $P_+$ truncates the moving-average kernel at lag zero, acting as the plus-part $\pluspart{\,\cdot\,}$ on the symbol. Composing with $N_+^{-1}$ gives \eqref{eq:filter}; for $\lambda>0$ the symbol $\hat n$ is bounded below, so $h$ and the position filter lie in $L^2$. The hypothesis $\int(1+\omega^2)S_\mu(\omega)/\hat n(\omega)\,d\omega<\infty$ places $h$ in a Sobolev class on which the half-line indicator multiplies boundedly \citep[Ch.~1]{LionsMagenes1972}, giving $\hat\chi\in L^2$. For the value, the first-order condition and adaptedness give $v=\tfrac12\E[x^\star_t\mu_t]$, and by the It\^o isometry with $\overline{\hat\psi}=\hat n_+\overline h$ on the real axis, $\E[x^\star\mu]=\tfrac1{2\pi}\|\pluspart{h}\|^2$, which is \eqref{eq:value}.

Referred to the rate, the same argument with $A=Q$ gives $\hat\chi=\hat q_+^{-1}\pluspart{h_\alpha}$ and $v=\frac1{4\pi}\|\pluspart{h_\alpha}\|^2$ with $h_\alpha=\hat q_-^{-1}\hat\varphi$, valid whenever $h_\alpha\in L^2$, the power law at $\lambda=0$ included. The two forms agree where both apply, since $\hat\psi=\varphi(0)+i\omega\hat\varphi$ makes $h-h_\alpha$ anticausal up to a lag-zero atom that the plus-part annihilates.

For the Ornstein--Uhlenbeck signal, $\hat\psi=\sigma\theta/(\theta-i\omega)$ and $h=\hat\psi\hat n_-^{-1}$; the combination
\[
h - \frac{\hat\psi}{\Phi(\theta)} = \sigma\theta\Bigl[\frac{1}{\hat n_-(\omega)} - \frac{1}{\hat n_-(-i\theta)}\Bigr]\frac{1}{\theta-i\omega}
\]
is anticausal --- the bracket's pole at $\omega=-i\theta$ cancels, $\hat n_-^{-1}$ is analytic in the lower half-plane, and $\hat n_-(-i\theta)=\hat n_+(i\theta)=\Phi(\theta)$ by conjugate symmetry. Hence $\pluspart{h}=\hat\psi/\Phi$, giving \eqref{eq:ou-filter} and, with $\|\hat\psi\|^2=\pi\sigma^2\theta$, the value $v=\sigma^2\theta/4\Phi^2$. For the power law at $\lambda=0$, $\hat q_-=(\gamma c_\beta)^{1/2}(-i\omega)^{-\nu}$ gives $\hat n_-(-i\theta)\propto\theta^{(1+\beta)/2}$ and, via \eqref{eq:vant}, the ratio $v/v_{\rm ant}=\sin(\pi\beta/2)$ of \eqref{eq:sinlaw}, the sine arising from the argument of $(-i\omega)^{\beta-1}$ on the two half-lines.

\section{Proof of Proposition~\ref{prop:boundary}}\label{app:C}
The finite-horizon policy \eqref{eq:finiteT} differs from the stationary one through the terminal-anchored weight and the truncation of the factor kernels to $[0,T]$. For the power-law factor \eqref{eq:gk-kernel} the truncation cuts the Marchaud tail at distance $d(t)$, contributing $d(t)^{-\nu}$; the weight deviation $((T-s)/(T-t))^\nu=1+O(d(t)^{-1})$ is subdominant. For a rational symbol the causal factor $\hat n_+^{-1}$ has kernel decaying as $e^{-b_1\tau}$ at the slowest zero $b_1$ of $\hat n_+$, so truncation at distance $d(t)$ costs $e^{-b_1 d(t)}$; equivalently the finite-horizon feedback gains relax to their stationary values at rate $2b_1$.

\paragraph{Acknowledgements.} The author thanks CFM colleagues, and in particular Jean-Yves Audibert, Antoine Bambade, Raphael Benichou, Jean-Philippe Bouchaud, Raphael Maillet, Iacopo Mastromatteo and Marc Potters, for useful discussions. Research and write-up were assisted by large language models. Verification scripts and figure code are available from the author.

\bibliographystyle{plainnat}
\bibliography{optimal-trading-filters}

\end{document}
